% !TeX root = document.tex
\subsection{Програмна реалізація двовимірної барицентричної інтерполяції}
У даному підрозділі розглядається еволюція програмної імплементації математичного апарату двовимірної барицентричної інтерполяції для задачі масштабування растрових зображень. Розробка відбувалася поступово: від створення базового прямого алгоритму, що слугував математичним еталоном, до його архітектурної та апаратної оптимізації. 

Кожен наступний етап розробки, який включав попереднє кешування вагових коефіцієнтів, застосування багатопотоковості, а також використання операцій лінійної алгебри, був спрямований на подолання критичних обчислювальних обмежень попередніх версій. Окремо розглядається перехід до обробки невеликих локальних матриць пікселів для усунення візуальних артефактів на різких градієнтах. 

Такий послідовний інженерний підхід дозволив перетворити теоретично ресурсомісткий алгоритм на високопродуктивне рішення, здатне працювати в режимі реального часу.
% !TeX root = document.tex
\subsubsection{Базова реалізація: прямий (нативний) метод}

Практична реалізація розробленого алгоритму починається з програмного представлення цифрового зображення як двовимірної дискретної матриці пікселів $F$ розмірністю $W \times H$. Граничні просторові індекси цієї матриці позначаються як $n$ та $m$, причому фізична ширина зображення становить $W = n + 1$, а висота — $H = m + 1$. Кожен елемент матриці з просторовими індексами $(i, j)$ зберігає значення кольору $f_{i,j}$ у форматі RGB. Задачею масштабування є побудова нової матриці зображення із цільовими розмірами $W_2 \times H_2$, де нова ширина $W_2 = N + 1$, а нова висота $H_2 = M + 1$ ($N$ та $M$ відповідають максимальним індексам цільової просторової сітки).

Першим етапом розробки стала програмна імплементація прямого (нативного) алгоритму двовимірної барицентричної інтерполяції, що представлений у коді функцією \texttt{processDirectBarycentricZoom}. Процес геометричного перетворення починається з відображення просторових кординат пікселів обох матриць у неперервний двовимірний інтервал $[-1, 1] \times [-1, 1]$. Для вихідного масиву даних просторова інтерполяційна сітка формується за законом розподілу екстремумів поліномів Чебишева другого роду:
\begin{equation}\label{eq:chebyshev_nodes}
	x_i = - \cos\left(\frac{\pi i}{n}\right), \quad y_j = - \cos\left(\frac{\pi j}{m}\right), \quad \text{де } i = \overline{0, n}, \quad j = \overline{0, m}.
\end{equation}

Так само просторові координати пікселів цільового зображення $(X, Y)$ переводяться у математичний діапазон $[-1, 1]$ за тим же правилом. У результаті ми отримуємо нормалізовані координати $(x, y)$, для яких безпосередньо обчислюється нове значення кольору:
\begin{equation}\label{eq:target_coordinates}
	x = - \cos\left(\frac{\pi X}{N}\right), \quad y = - \cos\left(\frac{\pi Y}{M}\right), \quad \text{де } X = \overline{0, N} \quad Y = \overline{0, M}.
\end{equation}

Для обчислення значення функції кольору $C(x, y)$ кожного результуючого пікселя використовується двовимірна друга барицентрична формула, яка враховує внесок усіх без винятку пікселів вхідного зображення:
\begin{equation}\label{eq:barycentric_2d_formula}
	C(x, y) = \frac{\sum\limits_{i=0}^{n} \vphantom{\sum}^{\prime} \sum\limits_{j=0}^{m} \vphantom{\sum}^{\prime} \frac{(-1)^{i+j}}{(x - x_i)(y - y_j)} f_{i,j}}{\sum\limits_{i=0}^{n} \vphantom{\sum}^{\prime} \sum\limits_{j=0}^{m} \vphantom{\sum}^{\prime} \frac{(-1)^{i+j}}{(x - x_i)(y - y_j)}}.
\end{equation}

З точки зору програмної архітектури, функція \texttt{processDirectBarycentricZoom} приймає об'єкт класу \texttt{QImage}, виконує попіксельне вилучення колірних компонент за допомогою низькорівневого доступу до пам'яті, ініціалізує результуючий об'єкт \texttt{QImage} заданих цільових розмірів і запускає обчислювальний процес. Важливою умовою стабільності програми є обробка ситуацій, коли поточні координати точки $(x, y)$ збігаються з вузлами вихідної просторової сітки Чебишева. У таких випадках знаменники у виразі \eqref{eq:barycentric_2d_formula} перетворюються на нуль, викликаючи критичні обчислювальні помилки. Для запобігання цьому у програмному коді реалізовано умовне розгалуження, логіка якого повністю спирається на математичне обґрунтування граничних випадків, що наведено вище. Якщо фіксується точний або частковий просторовий збіг, алгоритм автоматично блокує розрахунок загальної двовимірної форми та зводить операцію до одновимірного аналогу або безпосереднього копіювання значення кольору $f_{i,j}$.

Незважаючи на високу математичну точність, проведений аналіз обчислювальної ефективності прямого методу виявив його повну практичну непридатність для систем реального часу. Архітектура нативного алгоритму передбачає наявність чотирьох глибоко вкладених циклів: зовнішні цикли ітеруються по просторових координатах цільової матриці ($Y$ від $0$ до $M$, $X$ від $0$ до $N$), а всередині них для кожного кроку запускаються два вкладені цикли повного перебору вихідних точок ($j$ від $0$ до $m$, $i$ від $0$ до $n$). 

Така структура зумовлює вкрай високу асимптотичну складність алгоритму, яка оцінюється як $\mathcal{O}(W \cdot H \cdot W_2 \cdot H_2)$. Наприклад, якщо здійснюється просторове збільшення відносно невеликого графічного масиву розміром $500 \times 500$ пікселів у два рази (до розміру $1000 \times 1000$), процесору необхідно виконати:
\begin{equation}\label{eq:complexity_calc}
	1000 \times 1000 \times 500 \times 500 = 2.5 \cdot 10^{11}
\end{equation}
тобто 250 мільярдів ітерацій, кожна з яких містить операції ділення та множення. Таке колосальне навантаження викликає велике завантаження одного обчислювального потоку процесора, призводить до значних затримок (порядку кількох хвилин на один кадр) та спричиняє ефект тимчасового блокування (зависання) головного вікна графічного інтерфейсу користувача програми.

Реалізований прямий алгоритм на сітці Чебишева виконав свою основну роль у проекті — він підтвердив абсолютну коректность розробленого математичного апарату та виступив у ролі програмного еталона. Проте виявлені критичні обмеження швидкодії та нераціональне використання ресурсів процесора довели необхідність перебудови архітектури обчислень, що привело до розробки оптимізованої моделі.
\subsubsection{Оптимізація алгоритму: кешування ваг та багатопотокова обробка}

Фундаментальним недоліком прямого методу є багаторазове і надлишкове обчислення вагових коефіцієнтів у глибоко вкладених циклах. Оскільки ваги Чебишева залежать виключно від просторових координат вузлів сітки і не залежать від значень кольору самих пікселів, їх можна обчислити заздалегідь. Ця властивість стала основою для розробки оптимізованого алгоритму (\texttt{processBarycentricZoom}).

Першим етапом оптимізації є попереднє обчислення та збереження ваг. У програмному коді реалізовано спеціалізовану структуру даних \texttt{WeightCache}, яка зберігає масив вагових коефіцієнтів, їхню загальну суму (знаменник барицентричної формули), а також логічний прапорець \texttt{hit} та індекс \texttt{exact\_index}. Останні використовуються для миттєвої ідентифікації точного збігу координат цільової точки з вузлом вихідної сітки, що дозволяє безпечно уникати ділення на нуль без складних розгалужень у внутрішніх циклах обробки пікселів. Для максимального прискорення ініціалізації, розрахунок кеш-масивів для осей X та Y (\texttt{W\_X} та \texttt{W\_Y}) виконується паралельно у двох незалежних потоках (\texttt{std::thread}).

Другим етапом оптимізації стало усунення чотирьох вкладених циклів, які створювали надмірне обчислювальне навантаження. Замість них реалізовано двопрохідну схему обчислень із виділенням проміжного буфера пам'яті, що розділяє роботу алгоритму на два кроки:
\begin{enumerate}
	\item \textbf{Перший прохід (\texttt{pass1}):} Алгоритм ітерується по цільовій висоті $Y$ та вихідній ширині $i$, обчислюючи часткові барицентричні суми. Результати записуються у проміжний двовимірний масив \texttt{tempImg} розмірністю $H_2 \times (n+1)$. Для уникнення втрати точності на етапі проміжних обчислень, колірні компоненти RGB у \texttt{tempImg} зберігаються у форматі чисел з плаваючою комою (\texttt{float}).
	\item \textbf{Другий прохід (\texttt{pass2}):} Використовуючи попередньо підготовлені високоточні дані з \texttt{tempImg}, алгоритм виконує фінальне обчислення для кожної цільової координати $X$. Отримані значення нормалізуються, обмежуються стандартним колірним діапазоном $[0, 255]$ за допомогою функції \texttt{std::clamp} та записуються безпосередньо у результуючий об'єкт \texttt{QImage}.
\end{enumerate}

Ключову роль у забезпеченні швидкодії відіграє паралельна обробка даних на рівні центрального процесора. За допомогою визначення кількості доступних апаратних ядер (\texttt{std::thread::hardware\_concurrency}), виконання обох проходів розпаралелюється. Матриця цільового зображення розбивається на горизонтальні сегменти (chunks) вздовж осі Y. Кожен обчислювальний потік отримує свій незалежний блок рядків для обробки. Такий підхід гарантує, що потоки здійснюють запис у неперетинні області пам'яті, що повністю усуває необхідність використання механізмів блокування (наприклад, \texttt{std::mutex}) та запобігає виникненню стану гонитви (race condition).

Реалізація структурованого кешування у поєднанні з багатопотоковою двопрохідною архітектурою кардинально знизила обчислювальне навантаження. Замість екстремальної початкової складності $\mathcal{O}(W_2 \cdot H_2 \cdot W \cdot H)$, загальна кількість операцій тепер визначається сумою складностей двох послідовних проходів: $\mathcal{O}(H_2 \cdot W \cdot H)$ для першого етапу та $\mathcal{O}(H_2 \cdot W_2 \cdot W)$ для другого. Таким чином, загальна асимптотична складність оптимізованого алгоритму становить $\mathcal{O}(H_2 \cdot W \cdot (H + W_2))$. 

Для порівняння: при масштабуванні зображення розміром $500 \times 500$ пікселів у розмір $1000 \times 1000$, кількість ітерацій базового циклу зменшилася з 250 мільярдів (у прямому методі) до лише 750 мільйонів. Таке математичне скорочення кількості операцій більш ніж у 330 разів дозволило ефективно утилізувати апаратні ресурси багатоядерного процесора. Це перетворило алгоритм на практичне рішення, здатне обробляти зображення високої роздільної здатності за частки секунди без блокування графічного інтерфейсу користувача.
\subsubsection{Матрична форма двовимірної барицентричної інтерполяції}

Хоча класичний алгебраїчний запис двовимірної барицентричної формули~\eqref{eq:barycentric_2d_general} сам по собі забезпечує високу чисельну стійкість та алгоритмічну ефективність, наступним логічним етапом оптимізації розробленого методу є перехід до його матричного представлення.

Такий перехід зумовлений специфікою обробки цифрових зображень. Оскільки вагові коефіцієнти в барицентричній інтерполяції залежать виключно від просторових координат (вузлів сітки та цільових точок) і не залежать від значень інтенсивності пікселів, розрахунок геометричних ваг та власне інтерполяцію кольору доцільно розділити. Для забезпечення максимальної швидкодії при обробці зображень високої роздільної здатності математичну модель необхідно адаптувати під архітектуру сучасних обчислювальних систем.

У цьому підрозділі виконується математичний перехід від базового запису до матричної форми. Цей етап оптимізації дозволяє звести задачу двовимірного масштабування до двох послідовних матричних множень. З інженерної точки зору це є критично важливим кроком, оскільки матрична форма забезпечує кеш-дружній (послідовний) доступ до пам'яті, уможливлює розпаралелювання обчислень та відкриває шлях до глибокої апаратної оптимізації алгоритму за допомогою векторизованих інструкцій процесора.

\paragraph{Крок~1. Факторизація ядра.}
Розглянемо базову формулу для довільної цільової точки $(x, y)$:
\begin{equation}\label{eq:m_original}
	L_{n,m}(x, y) = \frac{
		\sideset{}{'}\sum\limits_{i=0}^{n}
		\sideset{}{'}\sum\limits_{j=0}^{m}
		\dfrac{(-1)^{i+j}}{(x - x_i)(y - y_j)}\, f_{i,j}
	}{
		\sideset{}{'}\sum\limits_{i=0}^{n}
		\sideset{}{'}\sum\limits_{j=0}^{m}
		\dfrac{(-1)^{i+j}}{(x - x_i)(y - y_j)}
	}.
\end{equation}

Введемо ненормалізовані вагові коефіцієнти, розділивши ядро на множники, що залежать лише від однієї змінної. Скористаємося тотожностями $(-1)^{i+j} = (-1)^i \cdot (-1)^j$ та $\frac{1}{(x-x_i)(y-y_j)} = \frac{1}{x-x_i} \cdot \frac{1}{y-y_j}$:
\begin{equation}\label{eq:m_uv}
	u_i(x) \stackrel{\text{def}}{=} \frac{w_i}{x - x_i}, \qquad
	v_j(y) \stackrel{\text{def}}{=} \frac{w_j}{y - y_j},
\end{equation}
де $w_k = (-1)^k$ з половинними значеннями на краях: $w_0 = \tfrac{1}{2}$, $w_n = \tfrac{(-1)^n}{2}$ (аналогічно для осі~$y$). Позначимо чисельник і знаменник формули~\eqref{eq:m_original} як $N$ та $D$ відповідно. Тоді:
\begin{equation}\label{eq:m_ND}
	N = \sum\limits_{i=0}^{n}\sum\limits_{j=0}^{m} u_i(x)\, v_j(y)\, f_{i,j}, \qquad
	D = \sum\limits_{i=0}^{n}\sum\limits_{j=0}^{m} u_i(x)\, v_j(y).
\end{equation}

\paragraph{Крок~2. Факторизація знаменника.}
Оскільки знаменник не містить множника $f_{i,j}$, що зв'язує індекси, подвійна сума розпадається у добуток двох одинарних:
\begin{equation}\label{eq:m_D_factor}
	D = \left(\sum\limits_{i=0}^{n} u_i(x)\right) \cdot \left(\sum\limits_{j=0}^{m} v_j(y)\right)
	\stackrel{\text{def}}{=} S_u(x) \cdot S_v(y).
\end{equation}

\paragraph{Крок~3. Розподіл ділення.}
Перегрупуємо чисельник - винесемо за внутрішню суму множники, що не залежать від~$j$:
\begin{equation}\label{eq:m_N_regroup}
	N = \sum\limits_{i=0}^{n} u_i(x) \cdot \left[\sum\limits_{j=0}^{m} v_j(y)\, f_{i,j}\right].
\end{equation}
Підставимо у вираз $L = N/D$:
\begin{equation}\label{eq:m_L_step}
	L = \frac{
		\sum\limits_{i=0}^{n} u_i(x) \cdot
		\left[\sum\limits_{j=0}^{m} v_j(y)\, f_{i,j}\right]
	}{S_u(x) \cdot S_v(y)}.
\end{equation}
Оскільки $S_v(y)$ не залежить від індексу~$i$, а $S_u(x)$ не залежить від індексу~$j$, можна розподілити ділення між обома сумами:
\begin{equation}\label{eq:m_L_distributed}
	L = \sum\limits_{i=0}^{n} \frac{u_i(x)}{S_u(x)} \cdot \sum\limits_{j=0}^{m} \frac{v_j(y)}{S_v(y)}\, f_{i,j}.
\end{equation}

\paragraph{Крок~4. Нормалізовані барицентричні коефіцієнти.}
Визначимо нормалізовані коефіцієнти:
\begin{equation}\label{eq:m_alpha_beta}
	\alpha_i(x) \stackrel{\text{def}}{=} \frac{u_i(x)}{S_u(x)} = \frac{\dfrac{w_i}{x - x_i}}{\sum\limits_{k=0}^{n} \dfrac{w_k}{x - x_k}}, \qquad
	\beta_j(y) \stackrel{\text{def}}{=} \frac{v_j(y)}{S_v(y)} = \frac{\dfrac{w_j}{y - y_j}}{\sum\limits_{k=0}^{m} \dfrac{w_k}{y - y_k}}.
\end{equation}
Безпосередньо з визначення випливає властивість нормалізації:
\begin{equation}\label{eq:m_norm_proof}
	\sum\limits_{i=0}^{n} \alpha_i(x) = \sum\limits_{i=0}^{n} \frac{u_i}{S_u} = \frac{1}{S_u} \sum\limits_{i=0}^{n} u_i = \frac{S_u}{S_u} = 1.
\end{equation}
Аналогічно $\sum_{j=0}^{m} \beta_j(y) = 1$. Іншими словами, $\alpha_i(x)$~--- це частка внеску $i$-го вузла у загальну суму ваг, і всі частки разом складають одиницю. Якщо $x$ збігається з вузлом $x_{i^*}$, покладаємо $\alpha_{i^*} = 1$, $\alpha_i = 0$ при $i \neq i^*$ (аналогічно для $\beta_j$).

У програмній реалізації нормалізація виконується одноразово на етапі
побудови вагових матриць: після обчислення всіх ваг
рядка знаходиться значення суми
$S_u(x_p)$, після чого кожен елемент рядка
ділиться на неї. Це повністю усуває операцію ділення з етапу
матричного множення.

\paragraph{Крок~5. Фінальна формула.}
Підставляючи~\eqref{eq:m_alpha_beta} у~\eqref{eq:m_L_distributed}, отримуємо:
\begin{equation}\label{eq:m_bilinear}
	L_{n,m}(x, y) = \sum\limits_{i=0}^{n}\sum\limits_{j=0}^{m} \alpha_i(x)\,\beta_j(y)\, f_{i,j}.
\end{equation}
Двовимірна барицентрична інтерполяція зведена до о зваженої суми значень $f_{i,j}$, де вага кожного пікселя є добутком двох одновимірних нормалізованих коефіцієнтів.

\paragraph{Крок~6. Матрична інтерпретація для однієї точки.}
Позначимо:
\begin{itemize}
	\item $\boldsymbol{\alpha}(x) = [\alpha_0(x),\; \alpha_1(x),\; \ldots,\; \alpha_n(x)]^{\!\top}$~--- вектор-стовпець розмірності $W$;
	\item $\boldsymbol{\beta}(y) = [\beta_0(y),\; \beta_1(y),\; \ldots,\; \beta_m(y)]^{\!\top}$~--- вектор-стовпець розмірності $H$;
	\item $\mathbf{F} \in \mathbb{R}^{H \times W}$~--- матриця кольорових значень вихідного зображення, $\mathbf{F}_{j,i} = f_{i,j}$.
\end{itemize}
Тоді формула~\eqref{eq:m_bilinear} записується як білінійна форма:
\begin{equation}\label{eq:m_bilinear_vec}
	L_{n,m}(x, y) = \boldsymbol{\beta}(y)^{\!\top}\,\mathbf{F}\, \boldsymbol{\alpha}(x),
\end{equation}
що обчислюється у два кроки:
\begin{enumerate}
	\item \textit{Інтерполяція вздовж $y$}:
	\begin{equation}\label{eq:m_pass1_vec}
		\mathbf{r}(y) = \mathbf{F}^{\!\top} \boldsymbol{\beta}(y) \in \mathbb{R}^{W}, \qquad
		r_i = \sum\limits_{j=0}^{m} \beta_j(y)\, f_{i,j}.
	\end{equation}
	\item \textit{Інтерполяція вздовж $x$}:
	\begin{equation}\label{eq:m_pass2_vec}
		L_{n,m}(x, y) = \boldsymbol{\alpha}(x)^{\!\top}\, \mathbf{r}(y) = \sum\limits_{i=0}^{n} \alpha_i(x)\, r_i.
	\end{equation}
\end{enumerate}

\paragraph{Крок~7. Узагальнення на всі цільові точки.}
Для цільового зображення розміром $W_2 \times H_2$ побудуємо дві \textit{вагові матриці}, рядки яких є нормалізованими барицентричними коефіцієнтами для кожного цільового пікселя:
\begin{equation}\label{eq:m_matrix_A}
	\mathbf{A} \in \mathbb{R}^{W_2 \times W}, \qquad \mathbf{A}_{p,i} = \alpha_i(x_p), \qquad p = 0, \ldots, W_2 - 1,
\end{equation}
\begin{equation}\label{eq:m_matrix_B}
	\mathbf{B} \in \mathbb{R}^{H_2 \times H}, \qquad \mathbf{B}_{q,j} = \beta_j(y_q), \qquad q = 0, \ldots, H_2 - 1,
\end{equation}
де $x_p = -\cos(\pi p / (W_2 - 1))$ та $y_q = -\cos(\pi q / (H_2 - 1))$. Кожен рядок обох матриць задовольняє умову нормалізації: сума елементів рядка дорівнює~$1$.

Елемент результуючої матриці $\mathbf{G}_{q,p} = L_{n,m}(x_p, y_q)$ відповідно до~\eqref{eq:m_bilinear} розраховується як:
\begin{equation*}
	\mathbf{G}_{q,p} = \sum\limits_{i=0}^{n}\sum\limits_{j=0}^{m} \mathbf{A}_{p,i}\, \mathbf{B}_{q,j}\, \mathbf{F}_{j,i}
	= \sum\limits_{i=0}^{n} \mathbf{A}_{p,i} \underbrace{\sum\limits_{j=0}^{m} \mathbf{B}_{q,j}\, \mathbf{F}_{j,i}} _{(\mathbf{B}\,\mathbf{F})_{q,i}}.
\end{equation*}
Внутрішня сума є елементом добутку матриць $\mathbf{B} \cdot \mathbf{F}$, а зовнішня~--- добутку з $\mathbf{A}^{\!\top}$. Отже, фінальне матричне рівняння має вигляд:
\begin{equation}\label{eq:m_result}
	\mathbf{G} = \mathbf{B}\, \mathbf{F}\, \mathbf{A}^{\!\top}.
\end{equation}

\paragraph{Поканальне представлення вхідних даних.}
Формула~\eqref{eq:m_result} визначена для матриці
$\mathbf{F}$. Оскільки колірне зображення формату RGB містить три
незалежні канали інтенсивності, у програмній реалізації матриця вхідних
даних декомпонується на три окремі матриці
$\mathbf{F}_R,\; \mathbf{F}_G,\; \mathbf{F}_B \in \mathbb{R}^{H \times W}$. Для кожного каналу матричне
множення виконується незалежно:
\begin{equation}\label{eq:m_result_rgb}
	\mathbf{G}_c = \mathbf{B}\, \mathbf{F}_c\, \mathbf{A}^{\!\top},
	\qquad c \in \{R, G, B\}.
\end{equation}
Такий підхід, на відміну від черезрядкового зберігання компонент, забезпечує послідовний доступ до пам'яті
всередині GEMM-ядра та сприяє автоматичній векторизації внутрішніх циклів
компілятором.

\paragraph{Програмна реалізація з бібліотекою Eigen.}
Для виконання матричного множення~\eqref{eq:m_result_rgb} використано
бібліотеку Eigen~--- header-only бібліотеку лінійної алгебри для C++, яка
реалізує апаратно оптимізований алгоритм GEMM
з автоматичним cache-blocking, використанням SIMD-інструкцій та
prefetch-директивами.

Вагові матриці $\mathbf{A}$ та $\mathbf{B}$ зберігаються у суцільних масивах типу \texttt{float} розмірністю $W_2 \times W$ та
$H_2 \times H$ відповідно. Для передачі їх до Eigen без копіювання даних
використовується механізм \texttt{Eigen::Map}, який створює матричний
об'єкт поверх існуючого масиву:
\begin{verbatim}
	Eigen::Map<MatRM> mWy(Wy.data(), H2, H);
	Eigen::Map<MatRM> mWx(Wx.data(), W2, W);
	Eigen::Map<MatRM> mFR(FR.data(), H, W);
\end{verbatim}
де \texttt{MatRM}~--- тип матриці у форматі
\texttt{Eigen::Matrix<float, Dynamic, Dynamic, RowMajor>}. Обчислення
результату для кожного каналу зводиться до єдиного виразу:
\begin{verbatim}
	MatRM mRR = (mWy * mFR) * mWx.transpose();
\end{verbatim}
Eigen автоматично створює проміжну матрицю
$\mathbf{T} = \mathbf{B} \cdot \mathbf{F}_c$ та виконує два послідовних
виклики оптимізованого GEMM-ядра.

Використання типу \texttt{float} замість \texttt{double} зменшує обсяг
даних удвічі, що подвоює ефективну пропускну здатність кеш-пам'яті та
дозволяє Eigen обробляти 8~елементів за один AVX-такт замість~4. Обчислення
самих вагових коефіцієнтів виконується у типі \texttt{double} для
збереження точності при діленні на малі різниці $(x - x_i)$, після чого
результат приводиться до \texttt{float}.

\paragraph{Побудова вагових матриць.}
Побудова матриць $\mathbf{A}$ та $\mathbf{B}$ виконується з використанням
усіх доступних апаратних потоків: діапазон цільових індексів рівномірно
розподіляється між потоками за допомогою пулу \texttt{std::thread}.
Спочатку всі доступні потоки будують матрицю $\mathbf{B}$, після завершення~- всі потоки будують
матрицю~$\mathbf{A}$.

На етапі побудови застосовано ті самі оптимізації, що й у злитому
алгоритмі: детекція точного збігу з вузлом за $O(1)$ на основі
цілочисельної арифметики, чергування знаків
та попередня нормалізація ваг, що усуває операцію ділення з етапу
GEMM-множення.

\paragraph{Конвертація вхідних даних.}
Вхідне зображення \texttt{QImage} конвертується у формат
\texttt{Format\_RGB32}, після чого колірні компоненти кожного пікселя
витягуються безпосередньо з \texttt{QRgb}-значень та записуються у три
окремі суцільні \texttt{float}-масиви $\mathbf{F}_R$, $\mathbf{F}_G$,
$\mathbf{F}_B$. Такий підхід усуває проміжні конвертації формату, що економить декілька повних проходів по масиву пікселів.

Результуючі значення після GEMM-множення обмежуються стандартним
колірним діапазоном $[0, 255]$ за допомогою \texttt{std::clamp} та
записуються безпосередньо у вихідний \texttt{QImage} формату
\texttt{Format\_RGB32} без проміжних матриць у цілочисельному форматі.

\paragraph{Обчислювальна складність.}
Побудова вагових матриць потребує $\mathcal{O}(W_2 \cdot W)$ операцій
для $\mathbf{A}$ та $\mathcal{O}(H_2 \cdot H)$ для~$\mathbf{B}$.
Матричне множення $\mathbf{B} \cdot \mathbf{F}_c$ виконується за
$\mathcal{O}(H_2 \cdot W \cdot H)$ операцій, а множення
$\mathbf{T}_c \cdot \mathbf{A}^{\!\top}$~--- за
$\mathcal{O}(H_2 \cdot W_2 \cdot W)$. Для трьох каналів загальна
асимптотична складність алгоритму становить:
\begin{equation}
	\mathcal{O}\bigl(3 \cdot W \cdot H_2 \cdot (H + W_2)\bigr)
	= \mathcal{O}\bigl(W \cdot H_2 \cdot (H + W_2)\bigr).
\end{equation}

Асимптотична складність матричної форми збігається зі складністю злитого
алгоритму. Однак матрична реалізація з використанням Eigen досягає
суттєво вищої практичної швидкодії завдяки оптимізованому GEMM-ядру:
cache-blocking забезпечує роботу з даними у кеш-пам'яті, SIMD-інструкції
обробляють 8~значень \texttt{float} за такт, а prefetch-директиви
зменшують затримку доступу до RAM. У сукупності ці фактори
забезпечують пропускну здатність GEMM, що наближається до теоретичної
пікової продуктивності процесора.
\subsubsection{Апаратна оптимізація: реалізація матричного методу засобами OpenCV}

Хоча базова матрична реалізація з використанням контейнерів \texttt{std::vector} та пулу потоків дозволила суттєво підвищити швидкодію, ручне управління операціями множення не дозволяє повною мірою використати апаратні можливості сучасних процесорів (наприклад, векторні SIMD-інструкції AVX2/AVX-512). Для досягнення максимальної продуктивності обчислень було прийнято рішення інтегрувати розроблений математичний апарат із бібліотекою комп'ютерного зору OpenCV (\texttt{processOpenCVBarycentricZoom}).

Головна перевага використання OpenCV полягає в наявності оптимізованого бекенду для операцій лінійної алгебри. Під капотом структури \texttt{cv::Mat} та її перевантажених операторів множення працюють високоефективні виклики до спеціалізованих математичних бібліотек (таких як BLAS, LAPACK або Eigen), які автоматично забезпечують апаратну векторизацію та оптимальне управління кешем на найнижчому рівні.

Алгоритмічний конвеєр OpenCV-реалізації складається з наступних кроків:
\begin{enumerate}
	\item \textbf{Конвертація типів:} Вихідний об'єкт \texttt{QImage} перетворюється на багатоканальну матрицю \texttt{cv::Mat}. Далі, за допомогою функції \texttt{cv::split}, зображення розшаровується на три окремі одноканальні матриці (канали R, G, B), що відповідають плоским масивам $F^{(c)}$ з попереднього методу.
	\item \textbf{Формування матриць ваг:} Нормалізовані вагові коефіцієнти Чебишева для обох осей генеруються за тими ж формулами \eqref{eq:matrix_Wy}, проте одразу записуються у структури \texttt{cv::Mat} типу \texttt{CV\_32F} (32-бітні числа з плаваючою комою). Це дозволяє уникнути накладних витрат на подальше копіювання пам'яті.
	\item \textbf{Оптимізоване множення:} Обчислення результуючої матриці кольору $C^{(c)}$ виконується за допомогою стандартного оператора множення матриць бібліотеки OpenCV, що математично еквівалентно виклику узагальненого множення матриць (GEMM - General Matrix Multiply):
	\begin{equation}\label{eq:opencv_gemm}
		C^{(c)} = W_y \times F^{(c)} \times W_x^T.
	\end{equation}
	Оскільки OpenCV автоматично розпаралелює ці операції (за наявності бекенду OpenMP або TBB), відпадає необхідність у ручному створенні пулу потоків \texttt{std::thread}.
	\item \textbf{Зворотне злиття:} Після незалежного обчислення нових розмірів для кожного колірного каналу, матриці об'єднуються назад у єдину багатоканальну структуру за допомогою функції \texttt{cv::merge}, нормалізуються і конвертуються назад у формат \texttt{QImage} для відображення у графічному інтерфейсі.
\end{enumerate}

Переведення математичної моделі на рейки індустріального стандарту OpenCV не змінило теоретичної асимптотичної складності алгоритму, яка залишилася на рівні $\mathcal{O}(H_2 \cdot W \cdot (H + W_2))$. Однак, завдяки глибокій оптимізації доступу до пам'яті та використанню векторних інструкцій процесора, практичний час виконання (wall-clock time) зменшився. Ця імплементація продемонструвала найкращий баланс між якістю зображення, притаманною барицентричній інтерполяції, та швидкістю обробки, наблизившись до показників вбудованих методів масштабування самої бібліотеки OpenCV.
\subsubsection{Кусково-поліноміальна інтерполяція: локальний барицентричний метод}

Глобальна барицентрична інтерполяція, розглянута у попередніх розділах, використовує значення всіх пікселів зображення для обчислення кольору кожної нової точки. З математичної точки зору це означає побудову інтерполяційного полінома надзвичайно високого ступеня (наприклад, ступеня 999 для рядка з 1000 пікселів). Хоча використання вузлів Чебишева нівелює феномен Рунге, на зображеннях із різкими просторовими переходами (наприклад, знімки екрана з текстом, графічні інтерфейси, термінали) глобальний підхід може провокувати появу артефактів дзвону (ringing artifacts, ефект Гіббса). Це проявляється у вигляді загасаючих хвилеподібних ліній (ехо-контурів), що поширюються від різких контрастних меж на сусідні області зображення.

Для усунення цієї проблеми було розроблено та програмно реалізовано алгоритм кусково-поліноміальної інтерполяції — локальний барицентричний метод (\texttt{processLocalBarycentric}). Його фундаментальна відмінність полягає у використанні концепції ковзного вікна (sliding window) фіксованого розміру $K \times K$, де $K$ — параметр згладжування (зазвичай 3, 4 або 6). 

Замість глобальної інтерполяції, значення цільового пікселя обчислюється виключно на основі його найближчого локального оточення. Алгоритмічний конвеєр локального методу складається з наступних етапів:
\begin{enumerate}
	\item \textbf{Проекція координат:} Цільові координати нормалізуються в діапазон $[-1, 1]$, після чого проектуються назад на неперервний простір вихідного зображення для знаходження умовної точки $(x_{src}, y_{src})$.
	\item \textbf{Визначення локального вікна:} Навколо знайденої точки обчислюються цілочисельні індекси початкового кута вікна $(x_{start}, y_{start})$. Розмір вікна жорстко обмежується краями вихідного зображення за допомогою функції \texttt{std::clamp} (або \texttt{std::min/max}) для уникнення виходу за межі пам'яті (out-of-bounds memory access).
	\item \textbf{Локальна нормалізація:} Координати умовної точки $(x_{src}, y_{src})$ перераховуються у локальну систему координат вікна $K \times K$, перетворюючись на локальні змінні $local\_x$ та $local\_y$ у діапазоні $[-1, 1]$.
	\item \textbf{Барицентрична оцінка:} Для обчислення фінального кольору застосовується класична двовимірна барицентрична формула \eqref{eq:barycentric_2d_formula}, але межі підсумовування обмежуються розміром вікна (від $0$ до $K-1$), а як значення $f_{i,j}$ використовуються виключно пікселі з поточного вікна.
\end{enumerate}

Локальні просторові вузли для вікна $K \times K$ генеруються за тим самим законом Чебишева-Лобатто \eqref{eq:chebyshev_nodes}, але адаптовані під кількість елементів $K$. Обчислення вагових коефіцієнтів та самої суми відбувається безпосередньо під час проходу по пікселях. Щоб уникнути ділення на нуль, алгоритм перевіряє потрапляння локальних координат $local\_x$ та $local\_y$ у вузли сітки (з точністю до $10^{-11}$). У разі точного збігу, колір відповідного пікселя з вихідного масиву копіюється у цільове зображення.

Використання локального вікна кардинально змінює обчислювальну складність алгоритму. Оскільки обчислення кольору кожної цільової точки вимагає перебору лише $K \times K$ сусідів, складність процесу масштабування стає повністю незалежною від розмірів вихідного зображення і становить $\mathcal{O}(H_2 \cdot W_2 \cdot K^2)$. 

Враховуючи, що $K$ є малою константою, локальний метод демонструє надзвичайно високу швидкість виконання, поступаючись лише вбудованим методам найближчого сусіда або білінійної інтерполяції. Програмна імплементація додатково оптимізована за рахунок багатопотоковості: матриця цільового зображення розбивається на горизонтальні сегменти, кожен з яких паралельно обробляється окремим потоком з пулу \texttt{std::thread}. 

Експериментальні дослідження підтвердили, що локальний барицентричний метод повністю усуває ефект дзвону на різких градієнтах, забезпечуючи ідеальне збереження контрастних меж при масштабуванні технічної графіки та тексту.